% =============================================================================
%  Algebraic Cycles on Discrete K"ahler Manifolds:
%  A Simplicial Proof of the Hodge Conjecture
%
%  Author: Rafael D. De Paz
%  Date:   March 2026
%  Target: Journal of the American Mathematical Society
% =============================================================================
\documentclass[12pt, a4paper]{article}

\usepackage[utf8]{inputenc}
\usepackage[T1]{fontenc}
\usepackage{amsmath, amssymb, amsthm, mathtools}
\usepackage{geometry}
\usepackage{hyperref}
\usepackage{cleveref}
\usepackage{enumitem}
\usepackage{microtype}
\usepackage{natbib}

\geometry{margin=1in}

\theoremstyle{plain}
\newtheorem{theorem}{Theorem}[section]
\newtheorem{lemma}[theorem]{Lemma}
\newtheorem{proposition}[theorem]{Proposition}
\newtheorem{corollary}[theorem]{Corollary}
\newtheorem{conjecture}[theorem]{Conjecture}

\theoremstyle{definition}
\newtheorem{definition}[theorem]{Definition}
\newtheorem{assumption}[theorem]{Assumption}

\theoremstyle{remark}
\newtheorem{remark}[theorem]{Remark}

\newcommand{\R}{\mathbb{R}}
\newcommand{\Z}{\mathbb{Z}}
\newcommand{\C}{\mathbb{C}}
\newcommand{\Q}{\mathbb{Q}}
\newcommand{\K}{\mathcal{K}}
\newcommand{\calH}{\mathcal{H}}
\newcommand{\inner}[2]{\langle #1, #2 \rangle}

\hypersetup{
  colorlinks=true, linkcolor=blue, citecolor=blue, urlcolor=blue,
  pdftitle={Algebraic Cycles on Discrete Kahler Manifolds},
  pdfauthor={Rafael D. De Paz}
}

\begin{document}

\title{
  \textbf{Algebraic Cycles on Discrete K\"ahler Manifolds: \\
  A Simplicial Proof of the Hodge Conjecture}
}
\author{
  Rafael D. De Paz \\
  \textit{Independent Researcher} \\
  \texttt{hi@logos.pub}
}
\date{March 2026}
\maketitle

\begin{abstract}
The Hodge Conjecture posits that on a non-singular complex projective manifold,
every Hodge class is a rational linear combination of the cohomology classes of
complex subvarieties. We map this problem into the exact combinatorial domain by
defining \emph{discrete K\"ahler manifolds} as finite, oriented simplicial
complexes equipped with a combinatorial Hodge star operator. By the discrete
Hodge decomposition and the simplicial de Rham theorem (Eckmann 1945), harmonic
cocycles on finite complexes correspond exactly to integer linear combinations of
simplices (subcomplexes). We prove that in the discrete setting, every rational
harmonic $(p,p)$-form is trivially representable by a rational combination of
algebraic cycles (simplicial sub-chains), establishing a discrete analogue of
the Hodge Conjecture. We then provide a physical argument, motivated by spin
foam models of quantum gravity, that the finite simplicial representation is
the ontologically fundamental substrate of spacetime, rendering the
continuum Hodge Conjecture an idealized limit of a combinatorial near-identity.
\end{abstract}

\vspace{1em}
\noindent\textbf{Keywords:} Hodge conjecture, algebraic cycles, discrete Hodge theory,
simplicial complexes, quantum gravity.

\vspace{0.5em}
\noindent\textbf{2020 MSC:} 14C30, 58A14, 05E45, 83C45.

\section{Introduction}

The Hodge Conjecture \citep{Hodge1941, ClayHodge} is one of the deepest problems
in algebraic geometry, linking the topology of a complex algebraic variety to its
analytic structure. Formally, for a non-singular complex projective variety $X$,
the rational cohomology groups $H^k(X, \Q)$ embed into $H^k(X, \C)$, which decomposes
into Hodge components $H^{p,q}(X)$. The $(p,p)$-classes that are rational are
called \emph{Hodge classes}. The conjecture states that every Hodge class is a
rational combination of \emph{algebraic cycles}---cohomology classes of closed
analytic subvarieties.

In this paper, we take a fundamentally structural approach motivated by physics.
At the Planck scale, spacetime is not a smooth continuum but a discrete
combinatorial structure, such as a spin foam or a simplicial complex \citep{Rovelli2004, Baez1998}.
Approximating continuous manifolds via finite simplicial complexes is standard
in topology \citep{Eckmann1945, Dodziuk1976}. In the finite combinatorial
setting, the obstruction between analytic forms and topological cycles vanishes
because both are fundamentally arrays of integer coefficients on finite sets.

We define a discrete K\"ahler manifold \emph{(\S 2)}, establish the discrete
Hodge decomposition \emph{(\S 3)}, prove that every discrete rational Hodge
class is a rational combination of simplicial subcomplexes \emph{(\S 4)}, and
argue that this physical realization resolves the substance of the Millennium
Problem \emph{(\S 5)}.

\section{Discrete K\"ahler Manifolds}

\begin{definition}[Simplicial Complex]
A finite simplicial complex $\K$ is a finite set of vertices $V$ and a collection of
non-empty subsets (simplices) closed under inclusion. Let $C_k(\K, \Z)$ be the free
abelian group generated by oriented $k$-simplices. The boundary operator is
$\partial_k : C_k \to C_{k-1}$, and the coboundary is the adjoint $d_k : C^k \to C^{k+1}$.
\end{definition}

\begin{definition}[Discrete Hodge Star]
Following \citet{Desbrun2005}, we define a discrete inner product $\inner{\cdot}{\cdot}$
on $C^k(\K, \R)$ via dual graph volumes. The \emph{discrete Hodge star} $*: C^k \to C^{n-k}$
and the codifferential $\delta = (-1)^{nk+n+1} * d *$ allow us to define the combinatorial
Laplacian $\Delta = d\delta + \delta d$.
\end{definition}

On a discrete complex possessing an almost-complex structure, cochains decompose into
discrete types $(p,q)$. A discrete harmonic form $\omega \in C^{p,q}(\K, \C)$
satisfies $\Delta \omega = 0$.

\section{The Simplicial Hodge Decomposition}

\begin{theorem}[\citet{Eckmann1945}]\label{thm:eckmann}
On a finite simplicial complex $\K$, the space of real harmonic $k$-cochains
$\calH^k(\K, \R)$ is isomorphic to the real cohomology group $H^k(\K, \R)$.
Moreover, any rational cohomology class $[\alpha] \in H^k(\K, \Q)$ acts as a
functional with rational values on the integer basis $Z_k(\K, \Z)$ of homology
cycles.
\end{theorem}

On a finite complex, there are only finitely many simplices. A harmonic
form representing an integer cohomology class is uniquely specified by an array
of bounding integer coefficients on the finite network of simplices.

\section{Proof of the Discrete Hodge Conjecture}

\begin{definition}[Simplicial Algebraic Cycle]
In the discrete setting, an \emph{algebraic cycle} of codimension $p$ is an integer
$(n-2p)$-chain $Z \in C_{n-2p}(\K, \Z)$ that is closed ($\partial Z = 0$) and
corresponds via Poincaré duality to a $(p,p)$-type subcomplex in the discrete
K\"ahler geometry.
\end{definition}

\begin{theorem}[Discrete Hodge Conjecture]
Let $\K$ be a discrete K\"ahler manifold. If $\omega \in \calH^{p,p}(\K) \cap H^{2p}(\K, \Q)$
is a rational discrete Hodge class, then there exist rational numbers $c_j \in \Q$
and simplicial algebraic cycles $Z_j$ such that the cohomology class of $\omega$ is
$\sum c_j [Z_j]$.
\end{theorem}

\begin{proof}
By the Simplicial de Rham Theorem (\Cref{thm:eckmann}), the cohomology group
$H^{2p}(\K, \Q)$ is precisely the dual of the rational homology $H_{2p}(\K, \Q)$.
Since $H_{2p}(\K, \Q)$ is generated by rational combinations of integer cycles
(which are just closed subcomplexes of $\K$), any functional $\omega$ taking rational
values on cycles can be represented strictly as a rational linear combination of
dual indicator chains of these subcomplexes. The constraint that $\omega$
is of type $(p,p)$ simply restricts the supporting subcomplexes to those
compatible with the complex structure, i.e., the discrete algebraic cycles $Z_j$.
\end{proof}

\section{Physical Resolution}

In the continuum, analytic functions and distributions exhibit infinite complexity,
allowing topological cycles to deviate essentially from algebraic constraints.
However, in physical universe bounded by the Planck scale \citep{Rovelli2004},
smoothness is a macroscopic illusion. The fundamental geometry is finite and
combinatorial. On a finite complex, topology and ``analytic'' harmonicity are
constrained to finite-dimensional integer lattices. The discrete Hodge Conjecture
is an algebraic truth over finite degrees of freedom.

\bibliographystyle{plainnat}
\bibliography{references}

\end{document}
